Toss a fair coin forever. The proportion of heads tends to one half, but how fast, how rarely does it stray far, and what
shape does the whole path take when you look from far away? This book answers with the great limit theorems of probability
in their sharp forms: Donsker's theorem, which turns every random walk into Brownian motion; Berry and Esseen's bound on the
error of the central limit theorem; the stable laws that replace the bell curve when variances are infinite; Stein's method
of approximation; Cram\'er's large deviations; the concentration of measure that makes high dimensions predictable; de
Finetti's theorem on exchangeable sequences; Birkhoff's ergodic theorem, which explains the digits of continued fractions;
and Furstenberg's double recurrence, which proves that a set of integers of positive density contains a three-term progression. It ends with the Poisson--Dirichlet law, which governs both the
cycles of a random permutation and the prime factors of a random integer, and with the cutoff that makes a deck of cards random after a handful of riffle shuffles. Quoted theorems are marked where they occur. The Selberg--Delange formula is the one deferred to Volume~15. Each chapter
opens with the story behind its ideas, ends with exercises and complete solutions.
